\chapter{IS 800:2007 Provisions Applied to Truss Members}
\label{app:is800}

This appendix states, derives where necessary, and discusses the
specific clauses of IS~800:2007~\cite{is800-2007} that the
constraint module \texttt{src/constraints/is800\_checks.py}
enforces on every truss candidate. The thesis proper
(\cref{sec:results:is800}) reports the numerical consequences of
these checks; this appendix documents the underlying code of practice
so a reviewer can verify formula-by-formula correspondence.

All formulae are written in SI units; truss benchmarks published
in imperial units (10-bar, 25-bar, 72-bar) are converted to SI
inside the caller before these provisions are evaluated.

\section{Scope, Material Constants, and Safety Factors}
\label{app:is800:scope}

IS~800:2007 is the Indian limit-state design code for general
construction in steel. The limit states considered for truss
members are:
\begin{itemize}
  \item \textbf{ultimate}: tension yielding (Clause 6.2), tension
        rupture (Clause 6.3), and compression buckling
        (Clause 7.1);
  \item \textbf{stability}: slenderness limits (Clause 3.8);
  \item \textbf{serviceability}: deflection limits
        (Clause 5.6.1).
\end{itemize}
The three remaining limit states that IS~800 lists
(fatigue, brittle fracture, fire) are not applicable to the
thesis benchmarks and are therefore not enforced.

\subsection{Default Material: Fe~410}
\label{app:is800:fe410}

The thesis uses \emph{structural steel grade Fe~410} throughout
(the default grade of IS~2062 for general construction). Its
characteristic strengths are:
\begin{equation}
  f_y = \SI{250}{MPa},
  \qquad
  f_u = \SI{410}{MPa},
  \qquad
  E = \SI{2.0e11}{Pa}.
\end{equation}
These values are hard-coded as \texttt{FY\_DEFAULT}/\texttt{FU\_DEFAULT}
in \texttt{is800\_checks.py} and can be over-ridden per-call by
the kwargs \texttt{fy}, \texttt{fu}.

\subsection{Partial Safety Factors}
\label{app:is800:gamma}

IS~800 Table~5 distinguishes two partial safety factors on
material resistance that are relevant to tension and compression
members:
\begin{equation}
  \gamma_{m0} = \num{1.10}
  \quad\text{(yielding at gross section)},
  \qquad
  \gamma_{m1} = \num{1.25}
  \quad\text{(ultimate at net section or in buckling)}.
\end{equation}
All reported resistances in this appendix are \emph{design}
resistances that already incorporate the relevant $\gamma_m$.
The optimiser treats member demand $N^{(e)}$ as the signed
factored axial force from the FEM and compares it directly against
these design resistances.

\subsection{Cross-Section Assumption: Solid Circular}
\label{app:is800:section}

Every bar in the thesis benchmarks is modelled as a solid circular
cross-section of area $A^{(e)}$. This reduces the number of
design variables from three (area + a shape descriptor pair) to
one (area alone) --- a modelling choice that matches the original
Kaveh benchmark conventions~\cite{kaveh2010} and makes reported
areas directly comparable with the published literature. For a
solid disc of area $A$ and radius $R=\sqrt{A/\pi}$,
the \emph{second moment of area} is $I=\pi R^4/4$ and the
\emph{radius of gyration} is:
\begin{equation}
  r_g \;=\; \sqrt{\frac{I}{A}} \;=\; \frac{R}{2}
  \;=\; \frac{1}{2}\sqrt{\frac{A}{\pi}}.
  \label{eq:is800:rg}
\end{equation}
Implementation: \texttt{radius\_of\_gyration\_circular} at
\texttt{is800\_checks.py:52}. A future extension to rolled tubes
or angles is a single-function override.

\section{Clause 3.8 --- Slenderness Limits}
\label{app:is800:clause3_8}

For any truss member the effective slenderness ratio is
\begin{equation}
  \lambda \;=\; \frac{K L}{r_g},
  \label{eq:is800:slenderness}
\end{equation}
where $K$ is the \emph{effective length factor} (theoretical pin-pin
ends give $K=\num{1.0}$, which is the value used throughout) and
$r_g$ is the radius of gyration from
\cref{eq:is800:rg}. IS~800 Table~3 places separate limits for
primary members on tension and compression:
\begin{equation}
  \lambda \le
  \begin{cases}
    \num{180} & \text{tension member}, \\
    \num{250} & \text{compression member}.
  \end{cases}
  \label{eq:is800:slend_limits}
\end{equation}
The tension limit is a constructibility / serviceability
requirement (to prevent sag under self-weight of long tension
ties); the compression limit is a stability requirement (long
compression members have residual imperfection amplitudes that
make design-office analysis unreliable above $\lambda\sim 250$).

Implementation: \texttt{check\_slenderness} at
\texttt{is800\_checks.py:65}. The constraint is enforced as
\emph{hard} (any violation marks the design infeasible); a
\num{1}-percent tolerance is applied as standard in the penalty
function to avoid knife-edge infeasibility on floating-point
evaluation.

\section{Clause 6.2 --- Tension Capacity, Gross-Section Yielding}
\label{app:is800:clause6_2}

A tension member first yields over its gross cross-sectional area
when the demand reaches
\begin{equation}
  T_{dg} \;=\; \frac{A_g\,f_y}{\gamma_{m0}}.
  \label{eq:is800:Tdg}
\end{equation}
$A_g$ is the gross area ($=A^{(e)}$ for the welded-end solid-bar
members used in the thesis). For Fe~410,
$T_{dg}/A_g = \SI{227.27}{MPa}$, a quantity that appears as a
horizontal line on every stress-versus-axial-force plot in
\cref{chap:results} where IS~800 is switched on.

For benchmarks without an IS~800 layer (the bare Kaveh-style
stress cap of the 10-, 25- and 72-bar problems) the admissible
tension stress is the standard $f_y=\SI{250}{MPa}$ applied
without the factor $1/\gamma_{m0}$. This difference accounts for
most of the weight penalty that IS~800 inclusion introduces in
\cref{sec:results:ablation:is800} --- not the compression
buckling curve, as one might first assume.

\section{Clause 6.3 --- Tension Capacity, Net-Section Rupture}
\label{app:is800:clause6_3}

Where a tension member is bolt-connected, the net area $A_n$ is
the gross area minus the bolt-hole deductions. IS~800 gives the
ultimate net-section rupture capacity as
\begin{equation}
  T_{dn} \;=\; \frac{0.9\,A_n\,f_u}{\gamma_{m1}}.
  \label{eq:is800:Tdn}
\end{equation}
The factor \num{0.9} is an empirical allowance for eccentricity
and shear lag at the end connection. For welded-end members
(the thesis default), $A_n = A_g$ and
\cref{eq:is800:Tdn} simplifies to
$T_{dn}=0.9 A_g f_u/\gamma_{m1}=\SI{295.2}{MPa}\cdot A_g$. This
exceeds $T_{dg}$ for Fe~410, so \cref{eq:is800:Tdg} is the
binding tension-capacity constraint for welded-end members. In
bolt-connected benchmarks it is necessary to check both.

Implementation: \texttt{check\_tension\_rupture} at
\texttt{is800\_checks.py:114}, kwarg
\texttt{net\_area\_ratio} defaults to \num{1.0} (welded ends).

\section{Clause 7.1 --- Compression Capacity (Perry-Robertson)}
\label{app:is800:clause7_1}

IS~800 Clause 7.1.2 gives the design compressive stress
$f_{cd}$ as a reduced yield strength
\begin{equation}
  f_{cd} \;=\; \chi\,\frac{f_y}{\gamma_{m0}},
  \qquad 0<\chi\le 1,
  \label{eq:is800:fcd}
\end{equation}
where the reduction factor $\chi$ is the closed-form
Perry-Robertson expression adopted from Eurocode~3:
\begin{equation}
  \chi \;=\; \frac{1}{\varphi + \sqrt{\varphi^2-\overline{\lambda}^2}},
  \qquad
  \varphi \;=\;
    \tfrac{1}{2}\bigl[\,
      1 + \alpha\bigl(\overline{\lambda}-\num{0.2}\bigr)
        + \overline{\lambda}^{2}
    \,\bigr],
  \label{eq:is800:chi}
\end{equation}
and the \emph{non-dimensional slenderness} $\overline{\lambda}$ is
the ratio of yield-stress to elastic-buckling-stress half-powered:
\begin{equation}
  \overline{\lambda}
  \;=\;\sqrt{\frac{f_y}{f_{cc}}},
  \qquad
  f_{cc} \;=\;\frac{\pi^2 E}{(KL/r_g)^2}.
  \label{eq:is800:lambda_bar}
\end{equation}
$f_{cc}$ is the Euler critical stress for an ideally straight
elastic pinned-pinned column; $\overline{\lambda}\to 0$ recovers
$\chi=1$ (unreduced yield), and $\overline{\lambda}\to\infty$
recovers $\chi=1/\overline{\lambda}^2$ (elastic Euler-buckling
limit).

The design compressive \emph{force} resistance is then
\begin{equation}
  P_d \;=\; A^{(e)}\,f_{cd}.
  \label{eq:is800:Pd}
\end{equation}
Implementation: \texttt{check\_compression} at
\texttt{is800\_checks.py:140}. In the penalty function,
$|N^{(e)}_{\text{compression}}|\le P_d$ is enforced with the
same \num{1}-percent tolerance as the tension check.

\subsection{Choice of Buckling Curve}
\label{app:is800:clause7_1:curve}

IS~800 Table~7 distinguishes four buckling-curve variants
$(\mathrm{a},\mathrm{b},\mathrm{c},\mathrm{d})$ via the
imperfection factor $\alpha$:
\begin{equation}
  \alpha \in \{\num{0.21},\num{0.34},\num{0.49},\num{0.76}\}
  \quad\text{for curves}\quad
  (\mathrm{a},\mathrm{b},\mathrm{c},\mathrm{d}).
\end{equation}
The curve to use depends on the cross-section geometry and the
buckling axis; IS~800 Table~7 is the authoritative mapping. For
the thesis solid-circular cross-section, Table~7 does not
explicitly enumerate an entry (the code tabulates rolled I, H,
angle, channel, tubular, and welded box sections). The
closest physical match is a rolled tubular section, which
Table~7 places on curve~\mbox{\emph{a}} ($\alpha=\num{0.21}$).
However, thesis results use curve~\mbox{\emph{b}}
($\alpha=\num{0.34}$) as a conservative middle-ground choice:
this matches the recommendation of
Subramanian~\cite{subramanian2008} for ``general-purpose
welded solid bars'' and errs on the safe side. Switching to
curve~\mbox{\emph{a}} is a one-line override
(\texttt{alpha=0.21}) in \texttt{check\_compression} and
typically lowers the optimised weight by $\num{1}$--$\num{3}\,\%$
on the benchmarks of
\cref{sec:results:ablation:is800}.

\subsection{Closed-Form Bounds on $\chi$}
\label{app:is800:chi_bounds}

Two trivial properties of \cref{eq:is800:chi} help debugging:
\begin{itemize}
  \item $\chi\le 1$ always, enforced explicitly in
    \texttt{is800\_checks.py:163} by a \texttt{min(chi,1.0)}.
    Without this clip, very-low-slenderness
    ($\overline{\lambda}\to 0$) cases can produce
    numerically-slightly-greater-than-unity $\chi$ due to
    round-off in $\sqrt{\varphi^2-\overline{\lambda}^2}$.
  \item $\chi > 0$ always, because $\varphi>\overline{\lambda}$
    for all $\overline{\lambda}>0$ and $\alpha>0$.
\end{itemize}
In the IS~800 study of
\cref{sec:results:is800}, the reported median $\chi$ values on
the 10-bar compression members are
$\chi\in[\num{0.42},\num{0.88}]$; a purely-yielding material
model would correspond to $\chi\equiv 1$, which is why
including IS~800 increases optimised weight.

\section{Clause 5.6.1 --- Serviceability Deflection Limit}
\label{app:is800:clause5_6_1}

IS~800 Clause 5.6.1 caps the maximum deflection of a structure
under its characteristic service load:
\begin{equation}
  \delta_{\max} \;\le\; \frac{L_{\mathrm{span}}}{325},
  \label{eq:is800:deflection}
\end{equation}
where the divisor \num{325} is the Table~6 default for
\emph{roof purlins and girts}; for typical rafters the divisor is
\num{240}, for cantilevers \num{150}, and for gantry girders
\num{750}. The thesis uses \num{325} as a generic default
because the benchmark loads are classical gravity loads and the
benchmark geometries are roof-truss-like. The divisor is a kwarg
(\texttt{divisor}) of \texttt{check\_deflection} and is over-ridable
per call without recompiling any higher-level code.

\paragraph{Observation.} For the specific Kaveh 10-bar geometry
($L_{\mathrm{span}}=\SI{720}{in}=\SI{18.3}{m}$), the IS~800
serviceability deflection limit at divisor~\num{325} is
$\delta_{\max}=\SI{56.3}{mm}$. The benchmark's original
Imperial-unit problem statement fixes the tip deflection limit
at \num{2.0}~in $=\SI{50.8}{mm}$, which is \emph{tighter} than
IS~800 serviceability would require. The 10-bar displacement
constraint therefore remains the binding serviceability check
in practice.

\section{Compliance Matrix: Code Checks per Benchmark}
\label{app:is800:matrix}

\Cref{tab:is800_matrix} summarises which of the above clauses is
enforced in each benchmark. Benchmarks whose original problem
statement already caps a stress or a deflection are kept in their
original units and limit unless the IS~800 ablation is
specifically switched on.

\begin{table}[h]
  \centering
  \caption{Per-benchmark IS~800 coverage. $\checkmark$ means the
    clause is enforced by default; ``--'' means the benchmark's
    original non-IS~800 limit is kept; ``opt.'' means IS~800 is
    enabled only for the ablation in
    \cref{sec:results:ablation:is800}.}
  \label{tab:is800_matrix}
  \begin{tabular}{l c c c c c}
    \toprule
    Clause / Benchmark & 10-bar & 25-bar & 72-bar & 200-bar & Notes \\
    \midrule
    3.8 slenderness       & opt. & opt. & opt. & \checkmark & 200-bar SI units \\
    6.2 tension yield     & opt. & opt. & opt. & \checkmark &  \\
    6.3 tension rupture   & opt. & opt. & opt. & \checkmark & welded ends \\
    7.1 compression       & opt. & opt. & opt. & \checkmark & curve~b default \\
    5.6.1 deflection      & --   & --   & --   & \checkmark & divisor~325 \\
    \bottomrule
  \end{tabular}
\end{table}

The 200-bar stepped-tower benchmark is the only one specified
natively in SI units and is therefore the only benchmark whose
default configuration has IS~800 fully engaged.

\section{Validation: Worked Numerical Check}
\label{app:is800:validation}

As a spot-check that \cref{eq:is800:chi}--\cref{eq:is800:Pd} is
implemented correctly, consider a single \num{2.0}-m solid-circular
bar of area $A=\SI{5e-4}{m^2}$ (diameter $\approx\SI{25.2}{mm}$)
under compression. Step-by-step:
\begin{align*}
  r_g &= \tfrac{1}{2}\sqrt{A/\pi}
       = \tfrac{1}{2}\sqrt{\SI{5e-4}{m^2}/\pi}
       = \SI{6.3e-3}{m}, \\
  \lambda &= KL/r_g = \SI{2.0}{m}/\SI{6.3e-3}{m} = \num{317.5}, \\
  f_{cc}  &= \pi^2 E / \lambda^2
          = \pi^2\cdot\SI{2.0e11}{Pa}/\num{317.5}^2
          = \SI{1.96e7}{Pa}, \\
  \overline{\lambda} &= \sqrt{f_y/f_{cc}}
          = \sqrt{\SI{2.5e8}{Pa}/\SI{1.96e7}{Pa}}
          = \num{3.57}, \\
  \varphi &= \tfrac{1}{2}\bigl[1+\num{0.34}(\num{3.57}-\num{0.2})+\num{3.57}^2\bigr]
          = \num{7.45}, \\
  \chi    &= \bigl[\num{7.45}+\sqrt{\num{7.45}^2-\num{3.57}^2}\bigr]^{-1}
          = \num{0.073}, \\
  f_{cd}  &= \chi f_y/\gamma_{m0}
          = \num{0.073}\cdot\SI{2.5e8}{Pa}/\num{1.1}
          = \SI{1.65e7}{Pa}, \\
  P_d     &= A\,f_{cd}
          = \SI{5e-4}{m^2}\cdot\SI{1.65e7}{Pa}
          = \SI{8.27e3}{N}.
\end{align*}
The slenderness is well above the \num{250}-limit, so the
Clause~3.8 check \emph{fails first}; a compliant cross-section
would need a larger radius of gyration. This is exactly the
behaviour observed on the 200-bar optimiser runs in
\cref{sec:results:200bar}, where small-area members at the tower
top are driven by the \num{250} compression-slenderness limit
rather than by the Perry-Robertson capacity.
